\section{Introduction}
\label{sec:intro}

% Motivation (1 paragraph): 3D printers are common, CAD skills are not. Text-to-3D mesh models make shapes that
% look right but have no exact dimensions, holes or wall thicknesses. Parametric CAD code is a better output for an
% LLM: exact, editable, and checkable by running it and measuring the result.
%\TODO{motivation: why text to CAD code, not text to mesh}

Multiple disciplines make use of 3D generated models: engineers, artists, and additive manufacturers --- just to name a few.
The industry standard representation of a 3D model is currently the STL format, as it's triangular meshes are easy for supporting software to work with.

Currently, the bar to entry for these disciplines is learning how to use CAD or other modelling software.
This issue has been a motivating factor behind agents that produce 3D representations of objects.
However, most of these models fall in the text-to-mesh realm, and do not allow for ease of measurement for specifications.
Instead, a text-to-code model is able to have deterministic checks ran against the code, allowing easy validation, repeatable results, and parametric exactness.

% Knowledge-systems framing (course focus). The agent can draw on four kinds of knowledge; the report asks which
% one helps which model:
%  1. retrieved examples from a hand-checked library (RAG),
%  2. feedback from running the code and checking the mesh (errors, measured size) -> repair loop,
%  3. knowledge in the weights: model size, and fine-tuning a small model on verified teacher builds,
%  4. an outside repository (Thingiverse) for what code can't build, versus sculpting with SDF code.

In addition, a text-to-code model allows more variety of language models to be used for orchestration (multi-modal models are not required).
For these orchestrators, there are multiple sources of knowledge that can be used in the aid of mesh structures:

\begin{itemize}
    \item Retrieved examples from a library (RAG)
    \item Code feedback
    \item Mesh bounding box checks
    \item Visual renders of meshes to check for accuracy
    \item Model fine-tuning
    \item Online services or repositories that host user-created models
\end{itemize}

% Research questions (each answered in Sec. 5 and in one line in the conclusion):
%  RQ1  How much does the model matter, and is the gap about size or running locally?        (5.1, 5.4)
%  RQ2  Do retrieved examples help, and for which models?                                     (5.2)
%  RQ3  Can fine-tuning on verified builds close the gap for a small local model?             (5.3)
%  RQ4  How much do the repair loop and retrying the whole request add?                       (5.5)
%  RQ5  For organic shapes, is sculpting with code better than downloading a model?          (5.6)
%  and, underneath all shape results: how far can a visual LLM judge be trusted?             (5.7)

This leads to the following research questions, which are addressed by this project:

\begin{enumerate}
    \item How much does the model matter, and is the gap about size or running locally?
    \item Do retrieved examples help, and for which models?
    \item Can fine-tuning on verified builds close the gap for a small local model?
    \item How much do the repair loop and retrying the whole request add?
    \item For organic shapes, is sculpting with code better than downloading a model?
    \item How far can a visual LLM judge be trusted? 
\end{enumerate}

% Contributions (4-5 bullets, each pointing to a section):
%  - An agent that turns a request into a checked STL: intake gate, RAG, CadQuery in a sandbox, mesh and size
%    checks, repair loop; SDF sculpting and Thingiverse for organic objects; combined parts (Sec. 3).
%  - A 48-case mechanical benchmark with hand-written reference solutions and expected sizes, plus 30 organic
%    cases (Sec. 4).
%  - A visual judge chosen among four on 146 blind hand grades, with the choice rule written before the runs
%    (Sec. 4.4, 5.7).
%  - One pre-planned round: 7 builders x 78 cases x 3 runs, case-level intervals and paired tests (Sec. 5).
%  - A fine-tuned 4B local model trained only on teacher builds that passed the checks (Sec. 3.6, 5.3).

In this project, we present an agent that turns a request into a verified STL.
We design a synthetic benchmark of 48 mechanical parts and 30 organic shapes.
An visual LLM judge was chosen based on a blind study of 146 user graded parts, which was used to evaluate the rest of the generations.
We also provide a training regimen for fine-tuning local models to better suit this task.
